\documentclass{article}
\usepackage{icml2025}

% ── Optional recommended packages ──────────────────────────────────────────
\usepackage[utf8]{inputenc}
\usepackage[T1]{fontenc}
\usepackage{amsmath,amssymb,amsthm}
\usepackage{mathtools}
\usepackage{booktabs}
\usepackage{multirow}
\usepackage{xcolor}
\usepackage{microtype}
\usepackage{enumitem}
\usepackage{algorithm}
\usepackage{algpseudocode}
\usepackage{graphicx}
\usepackage{tikz}
\usetikzlibrary{arrows.meta,positioning,shapes.geometric,fit,backgrounds}

% ── Theorem environments ───────────────────────────────────────────────────
\newtheorem{theorem}{Theorem}
\newtheorem{proposition}[theorem]{Proposition}
\newtheorem{corollary}[theorem]{Corollary}
\newtheorem{definition}{Definition}
\newtheorem{remark}{Remark}

% ── Notation shorthands ────────────────────────────────────────────────────
\newcommand{\Prob}{\mathbb{P}}
\newcommand{\Nat}{\mathbb{N}}
\newcommand{\Real}{\mathbb{R}}
\newcommand{\calS}{\mathcal{S}}
\newcommand{\calP}{\mathcal{P}}
\newcommand{\Distr}{\mathrm{Distr}}
\newcommand{\topka}{T\!O\!P^\alpha_k}
\newcommand{\eps}{\varepsilon}
\newcommand{\bigO}{\mathcal{O}}

\icmltitlerunning{GPU-Parallel PCTL Verification of LLM Outputs}

% =============================================================================
\begin{document}

\twocolumn[
\icmltitle{GPU-Parallel Probabilistic Model Checking of\\
           Large Language Model Outputs}

\begin{icmlauthorlist}
\icmlauthor{Anonymous Author(s)}{anon}
\end{icmlauthorlist}
\icmlaffiliation{anon}{Anonymised for review}

\icmlkeywords{probabilistic model checking, large language models, PCTL,
              statistical model checking, GPU inference}

\vskip 0.3in
]

\printAffiliationsAndNotice{\icmlanonymize}

% ── Abstract ──────────────────────────────────────────────────────────────
\begin{abstract}
Probabilistic Computation Tree Logic (PCTL) model checking of Large Language
Model (LLM) text generation provides formal, probability-certified answers to
questions such as \emph{``What is the probability this model generates
gender-biased text in the next $L$ tokens?''}
The recently proposed \textsc{LLMChecker} framework~\citep{gross2025llmchecker}
demonstrates this is feasible for small lookahead depths, but is bottlenecked by
state-space explosion: the $\alpha$-$k$-bounded DTMC has up to $k^L$ states,
and exact verification via the Storm model checker becomes intractable beyond
$L \approx 10$ tokens.
We make three contributions.
\textbf{(1)} We replace the recursive depth-first DTMC construction with a
\emph{breadth-first} construction that submits all nodes at depth $d$ as a
single batched GPU inference call, reducing serial LLM calls from $\bigO(k^L)$
to $\bigO(L)$.
\textbf{(2)} We exploit the tree structure of the DTMC to replace Storm's
iterative solver with a single GPU sparse backward induction pass
($\bigO(L)$ sparse matrix-vector products), reducing verification from hours
to milliseconds.
\textbf{(3)} We introduce \emph{Direct Statistical Model Checking} (Direct
SMC), which bypasses DTMC construction entirely by sampling $M$ complete
trajectories directly from the LLM, then evaluating the PCTL property on each.
Direct SMC scales to any $L$ bounded only by the model's context window.
We evaluate across the full GPT-2 family, seven PCTL property classes, a
systematic exact-vs-SMC boundary grid, a 30-prompt diversity study, and a
temperature sensitivity study — demonstrating $L=50$ token certificates in
under 16 seconds on a consumer GPU.
\end{abstract}

% =============================================================================
\section{Introduction}
\label{sec:intro}

Large Language Models generate text by sampling tokens sequentially from a
learned conditional distribution.
This sequential, probabilistic process has a natural mathematical
representation: a \emph{Discrete-Time Markov Chain} (DTMC) whose states are
partially generated strings and whose transitions are token probabilities.
\emph{Probabilistic Computation Tree Logic} (PCTL)~\citep{hansson1994logic,baier2008principles}
is a temporal logic for DTMCs, offering a vocabulary rich enough to express
diverse safety and quality properties of LLM outputs. For example:
\begin{itemize}[leftmargin=*,itemsep=1pt,topsep=2pt]
  \item $\Prob(\mathbf{F}\;\textit{gender}{>}0)$: probability of eventually
        generating male-biased text
  \item $\Prob(\mathbf{G}\;\textit{readability}{>}5997)$: probability of
        always producing readable text
  \item $\Prob(\mathbf{F}\;\textit{step}{=}5 \wedge \textit{sim}{>}90)$:
        probability of reproducing a copyrighted phrase
\end{itemize}
These are formal, probability-certified statements about the model's generation
mechanism — analogous to hardware circuit verification or safety proofs in
reactive systems.

\paragraph{The \textsc{LLMChecker} paper.}
\citet{gross2025llmchecker} introduced \textsc{LLMChecker}, the first system to
apply PCTL model checking to LLM text generation.
Their key idea is the $\alpha$-$k$ bounding trick: retain only the top-$k$
tokens whose cumulative probability reaches threshold $\alpha$.
Empirically, LLMs concentrate most mass on very few tokens (e.g.\ Gemma-2B
reaches $\alpha{=}0.9$ with the top 3 tokens), so $k$ can be kept small.
The resulting DTMC is encoded in PRISM and verified by
Storm~\citep{hensel2022storm}.

\paragraph{The state-space explosion problem.}
The DTMC still has up to $\Theta(k^L)$ states.
For Mistral-7B with $k{=}15$, $L{=}5$: $255{,}085$ states and $23{,}023$
seconds of encoding time~\citep{gross2025llmchecker}.
For $L{=}30$, $k{=}5$: $5^{30} \approx 9.3 \times 10^{20}$ states — beyond
any computer.

\paragraph{Our contributions.}
We identify two unexploited properties of the \textsc{LLMChecker} DTMC:
\begin{enumerate}
  \item \textbf{Level independence.}  All nodes at depth $d$ are mutually
    independent and require identical operations, enabling full GPU parallelism.
  \item \textbf{Tree structure.}  Because each DTMC state is a unique string
    prefix, there are no back-edges.  PCTL reachability on a tree
    reduces to simple backward induction, not iterative linear system solving.
\end{enumerate}
Exploiting (1) gives our BFS construction (\S\ref{sec:bfs}).
Exploiting (2) gives our GPU sparse backward induction (\S\ref{sec:backward}).
When $k^L$ itself is the barrier, we introduce Direct SMC (\S\ref{sec:dsmc}).

% =============================================================================
\section{Background}
\label{sec:background}

\subsection{Discrete-Time Markov Chains and PCTL}

\begin{definition}[DTMC]
A \emph{discrete-time Markov chain} is a tuple
$(S, P, s_0, \mathit{AP}, \mathit{Lab})$ where $S$ is a finite set of states,
$P: S \times S \to [0,1]$ with $\sum_{s'} P(s,s') = 1$ for all $s$,
$s_0 \in S$ is the initial state, and $\mathit{Lab}: S \to 2^{\mathit{AP}}$
labels states with atomic propositions.
\end{definition}

PCTL~\citep{hansson1994logic} extends CTL with a probabilistic operator.
We focus on the two path formulas from~\citet{gross2025llmchecker}:
\begin{align}
  \mathcal{P}(\mathbf{F}\,\phi)
    &\;\text{: probability of \emph{eventually} reaching $\phi$} \label{eq:eventually}\\
  \mathcal{P}(\mathbf{G}\,\phi)
    &\;\text{: probability of \emph{always} staying in $\phi$} \label{eq:always}
\end{align}
where $\phi$ is an atomic proposition over state features
(e.g.\ $\textit{gender} > 0$).

\subsection{LLM Formalisation and $\alpha$-$k$ Bounding}

Following~\citet{gross2025llmchecker}, an LLM is a function
$f: \Sigma^* \to \Distr(\Sigma)$ mapping a string to a probability distribution
over tokens.
The $\alpha$-$k$ selection operator retains the most probable tokens up to
cumulative probability $\alpha$, capped at $k$:
\begin{multline}
  \topka(f(\omega)) = \{(\sigma_i, p_i)\}_{i=1}^{n}, \\
  n = \min\!\bigl(k,\;\min\{m : \textstyle\sum_{i=1}^m p_i \ge \alpha\}\bigr).
\end{multline}
The $\alpha$-$k$-bounded DTMC has one state per distinct string prefix and one
terminal \textsc{Rest} state per node absorbing the remaining probability mass.

\subsection{The Original \textsc{LLMChecker} Algorithm}

Algorithm~\ref{alg:original} reproduces the DFS encoding procedure from
\citet{gross2025llmchecker}.
It issues one LLM call per node — $\sum_{d=0}^{L} k^d = \Theta(k^L)$ serial
calls in total, taking up to $23{,}023$ seconds for a single query.

\begin{algorithm}[tb]
\caption{Original DFS DTMC construction~\citep{gross2025llmchecker}}
\label{alg:original}
\begin{algorithmic}[1]
\Require $\omega, k, L, M(\cdot)$
\State $T \gets \topka(f(\omega))$ \Comment{\textbf{Serial} LLM call per node}
\For{$(\sigma, p) \in T$}
  \State create child $(\omega\cdot\sigma,\, p)$; recurse with $L{-}1$
\EndFor
\State create REST child with prob $1 - \sum p_i$; recurse
\end{algorithmic}
\end{algorithm}

% =============================================================================
\section{BFS Construction with Batched GPU Inference}
\label{sec:bfs}

\begin{proposition}[Level independence]
\label{prop:level-independence}
All nodes at depth $d$ are mutually independent: they share no state,
require identical operations ($\topka(f(\cdot))$), and their string prefixes
are fully determined before depth $d$ is processed.
\end{proposition}

\begin{proof}
Each node's string $\omega$ uniquely identifies its position.
The DTMC is a directed tree, so no transition leads from depth $d$ back to an
earlier depth; nodes at depth $d$ are thus causally independent.
\end{proof}

Proposition~\ref{prop:level-independence} motivates Algorithm~\ref{alg:bfs}:
collect all nodes at depth $d$, submit them as one batched GPU inference call,
and receive all top-$k$ distributions simultaneously.

\begin{algorithm}[tb]
\caption{BFS DTMC construction with batched GPU inference}
\label{alg:bfs}
\begin{algorithmic}[1]
\Require $\omega^{(0)}, L, \alpha, k, M(\cdot)$, backend
\Ensure level lists $\ell_0, \ldots, \ell_L$
\State active $\gets$ [root$(\omega^{(0)})$]
\For{$d = 0, \ldots, L{-}1$}
  \State res $\gets$ backend.get\_top\_k(active, $\alpha$, $k$) \label{line:batch}
    \Comment{one GPU call, all of level $d$}
  \For{(node, tokens, probs) $\in$ zip(active, res)}
    \State create child nodes; add REST absorbing child
  \EndFor
  \State active $\gets$ non-REST children
\EndFor
\end{algorithmic}
\end{algorithm}

\begin{proposition}[Serial LLM calls]
Algorithm~\ref{alg:original} requires $\Theta(k^L)$ serial LLM calls.
Algorithm~\ref{alg:bfs} requires exactly $L$ batched LLM calls.
\end{proposition}

\paragraph{vLLM prefix caching synergy.}
When paired with vLLM~\citep{kwon2023efficient}'s automatic KV-cache prefix
sharing, all $k^d$ nodes at depth $d$ share the common prefix $\omega^{(0)}$.
vLLM's PagedAttention stores KV blocks in a content-addressed cache, serving
already-computed prefixes at zero additional cost.
We observe two orders-of-magnitude throughput growth from depth 0 (12K tok/s,
cold cache) to depth 3 (48M tok/s, near-full cache hits) during BFS
construction, confirming near-complete prefix reuse.

% =============================================================================
\section{GPU Sparse Backward Induction}
\label{sec:backward}

\begin{proposition}[DTMC is a tree]
The $\alpha$-$k$-bounded DTMC produced by Algorithm~\ref{alg:bfs} is a directed
tree: each state is a unique string prefix, so there are no cycles.
\end{proposition}

\begin{corollary}
PCTL reachability on a tree DTMC reduces to finite backward induction; no
iterative linear system solver is required.
\end{corollary}

Let $V_d[i]$ denote the PCTL probability at node $i$ of depth $d$.

\paragraph{Eventually $\Prob(\mathbf{F}\,\phi)$.}
\begin{equation}
  V_d[i] =
  \begin{cases}
    1 & \phi(\text{node}_i) \\
    0 & \text{node}_i\text{ terminal} \\
    \sum_j T_d[i,j]\cdot V_{d+1}[j] & \text{otherwise}
  \end{cases}
  \label{eq:eventually-recurrence}
\end{equation}

\paragraph{Always $\Prob(\mathbf{G}\,\phi)$.}
\begin{equation}
  V_d[i] =
  \begin{cases}
    0 & \neg\phi(\text{node}_i)\;\text{or terminal} \\
    \sum_j T_d[i,j]\cdot V_{d+1}[j] & \text{otherwise}
  \end{cases}
  \label{eq:always-recurrence}
\end{equation}

where $T_d \in \Real^{|\ell_d| \times |\ell_{d+1}|}$ is the sparse transition
matrix at depth $d$.
Each backward step is a single \texttt{torch.sparse.mm} call on GPU.
The total number of GPU kernel launches is $L$, compared to Storm's iterative
solver which requires $O(|S|)$ matrix-vector products.
In practice verification times drop from hours (Storm, CPU) to milliseconds.

\begin{algorithm}[tb]
\caption{GPU sparse backward induction}
\label{alg:bwd}
\begin{algorithmic}[1]
\Require levels $\ell_0,\ldots,\ell_L$; operator ($\mathbf{F}$/$\mathbf{G}$); predicate $\phi$
\Ensure $V_0[0] = \Prob(\mathbf{F}/\mathbf{G}\;\phi)$
\State $V_L[i] \gets \phi(\text{node}_i)$ for all $i \in \ell_L$
\For{$d = L{-}1$ \textbf{downto} $0$}
  \State Build sparse $T_d$ (COO) from $\ell_d$ children
  \State $V_d \gets T_d \cdot V_{d+1}$ \Comment{\texttt{torch.sparse.mm}}
  \State Apply $\mathbf{F}$/$\mathbf{G}$ boundary conditions
\EndFor
\Return $V_0[0]$
\end{algorithmic}
\end{algorithm}

% =============================================================================
\section{Direct SMC for Long Lookaheads}
\label{sec:dsmc}

\subsection{Motivation}

Even with BFS and GPU backward induction, the $\Theta(k^L)$ state space is the
fundamental barrier.
For $L{=}50$, $k{=}5$: $5^{50} \approx 8.9 \times 10^{34}$ states — no
algorithm can enumerate these.
PCTL reachability probabilities can be \emph{estimated} without enumeration
using Statistical Model Checking (SMC)~\citep{younes2002probabilistic,legay2010statistical}.
We introduce \emph{Direct SMC}: rather than sampling paths through an explicit
DTMC (which still requires building the tree), we sample directly from the
LLM's generation process.

\subsection{Algorithm}

\begin{algorithm}[tb]
\caption{Direct SMC for long-lookahead PCTL}
\label{alg:dsmc}
\begin{algorithmic}[1]
\Require prompt $\omega^{(0)}$, depth $L$, predicate $\phi$, quantifier $M$,
         sample count $N$, confidence $\delta$
\Ensure estimate $\hat{p}$, CI $[\hat{p}{-}\eps,\,\hat{p}{+}\eps]$
\State count $\gets 0$
\For{each chunk of $B$ trajectories until $N$ total}
  \State generate $B$ trajectories of length $L$ \Comment{one GPU call}
  \State increment count for each satisfied trajectory
\EndFor
\State $\hat{p} \gets \text{count}/N$;\quad
       $\eps \gets \sqrt{\ln(2/\delta)/(2N)}$
\Return $\hat{p},\;\max(0,\hat{p}{-}\eps),\;\min(1,\hat{p}{+}\eps)$
\end{algorithmic}
\end{algorithm}

Let $\omega_t = \omega^{(0)}\sigma_1\cdots\sigma_t$.
\textsc{Evaluate} checks satisfaction along a trajectory:
\begin{align}
  \text{sat}_{\mathbf{F}} &= \textstyle\bigvee_{t=1}^{L} \phi(M(\omega_t), t) \label{eq:sat-F}\\
  \text{sat}_{\mathbf{G}} &= \textstyle\bigwedge_{t=1}^{L} \phi(M(\omega_t), t) \label{eq:sat-G}
\end{align}

\subsection{Statistical Guarantee}

\begin{theorem}[Chernoff--Hoeffding bound]
\label{thm:chernoff}
Let $p$ be the true PCTL probability and $\hat{p} = \text{count}/M$ the Direct
SMC estimate from $M$ i.i.d.\ trajectory samples.
For any $\delta \in (0,1)$:
\[
  \Prob\!\left(|\hat{p} - p| \le \sqrt{\frac{\ln(2/\delta)}{2M}}\right) \ge 1 - \delta.
\]
\end{theorem}

\begin{proof}
Each trajectory evaluation is an independent Bernoulli trial.
Hoeffding's inequality~\citep{hoeffding1963} applied to the bounded $[0,1]$
random variable gives the result.
\end{proof}

\begin{corollary}[Sample complexity]
For additive error $\eps{=}0.025$ at confidence $1{-}\delta{=}0.95$:
$M \ge \ln(2/\delta)/(2\eps^2) = 3{,}689$ samples suffice.
\end{corollary}

\paragraph{vLLM integration.}
\texttt{SamplingParams(n=}$B$\texttt{)} generates $B$ independent completions
in one forward-pass session.
All $B$ trajectories share the KV cache for $\omega^{(0)}$ (computed once),
amortising the prompt prefill cost.

\subsection{Complexity Comparison}

\begin{table}[tb]
\centering\small
\caption{Method comparison. ``Serial'' = sequential LLM calls;
``Total'' = all LLM calls; ctx = context window length.}
\label{tab:complexity}
\begin{tabular}{lccc}
\toprule
Method & Serial & Total & Max $L$ \\
\midrule
DFS exact~\citep{gross2025llmchecker} & $\Theta(k^L)$ & $\Theta(k^L)$ & ${\approx}10$ \\
BFS exact (ours) & $L$ & $\Theta(k^L)$ & ${\approx}8$ \\
Direct SMC (ours) & $M/B$ & $M{\cdot}L$ & ${\le}$ctx \\
\bottomrule
\end{tabular}
\end{table}

% =============================================================================
\section{Text Quantification and PCTL Vocabulary}
\label{sec:quant}

PCTL atomic propositions reference integer-valued features of the generated
text, computed by a \emph{quantification function} $M: \Sigma^* \times \Nat \to
\mathbb{Z}^n$.
Table~\ref{tab:quant} lists the original methods from~\citet{gross2025llmchecker}
and five new quantifier classes we introduce.

\begin{table}[tb]
\centering
\small
\caption{Text quantification methods. Top block: original \textsc{LLMChecker}
features. Lower blocks: new contributions.}
\label{tab:quant}
\begin{tabular}{llp{2.5cm}}
\toprule
Feature & Range & Example predicate \\
\midrule
\multicolumn{3}{l}{\textit{Original quantifiers}} \\
\textit{gender} & $\mathbb{Z}$ & $\textit{gender} > 0$ \\
\textit{polarity} & $[-100, 100]$ & $\textit{polarity} \ge 30$ \\
\textit{readability} & $[0, 10^4]$ & $\textit{readability} > 5997$ \\
\textit{similarity} & $[0, 100]$ & $\textit{similarity} > 90$ \\
\textit{step} & $\Nat$ & $\textit{step} = 5$ \\
\midrule
\multicolumn{3}{l}{\textit{Lexical / symbolic (new)}} \\
\textit{keywords} & $\Nat$ & $\textit{keywords} > 0$ \\
\textit{forbidden} & $\Nat$ & $\textit{forbidden} = 0$ \\
\textit{json} & $\{0, 100\}$ & $\textit{json} \ge 100$ \\
\midrule
\multicolumn{3}{l}{\textit{Classifier-backed (new)}} \\
\textit{toxicity} & $[0, 100]$ & $\textit{toxicity} < 50$ \\
\bottomrule
\end{tabular}
\end{table}

\paragraph{Lexical and symbolic quantifiers.}
\textbf{KeywordPresence} counts keyword hits, supporting crisp reachability
queries (\emph{``does the model produce an answer marker in $L$ tokens?''}).
\textbf{ForbiddenWordChecker} paired with $\Prob(\mathbf{G}\;\textit{forbidden}{=}0)$
yields a formal probability certificate that the model \emph{never} produces a
disallowed token.
\textbf{FormatConstraint} supports \texttt{json} (parseable JSON in output) and
\texttt{balanced\_parens} modes, making structural properties first-class PCTL
citizens.

\paragraph{Classifier-backed quantifiers.}
\textbf{ToxicityHeuristic} uses a curated word list as a lightweight baseline.
\textbf{ClassifierQuantifier} wraps any HuggingFace \texttt{text-classification}
pipeline (e.g.\ \texttt{unitary/toxic-bert}):
\[
  \Prob(\mathbf{G}\;\textit{toxicity} < 50)
  \quad\text{``model stays below 50\% toxicity throughout generation''}
\]
This is the single largest scope multiplier: it moves from lexical counting
toward semantically grounded safety certificates with no changes to the
verification infrastructure.

\paragraph{Conjunctive timed properties.}
Pairing \textit{step} with any feature enables timed specifications:
\[
  \Prob\bigl(\mathbf{F}\;\textit{step}{=}10 \wedge \textit{gender}{>}0\bigr)
  \quad\text{``male-biased text by exactly token 10''}
\]

% =============================================================================
\section{Experiments}
\label{sec:experiments}

\subsection{Setup}

All experiments use GPT-2 (124M) at temperature 1.0 unless otherwise stated.
Section~\ref{sec:model-scaling} additionally tests GPT-2-medium (355M), -large
(774M), and -XL (1.5B).
Hardware: NVIDIA RTX~5090 (32\,GB VRAM), Python 3.11, PyTorch 2.10,
HuggingFace Transformers.
The original \textsc{LLMChecker} experiments used dual AMD Epyc 7763 CPUs with
A100 GPUs; our results show the GPU-parallel approach is accessible on consumer
hardware.

\subsection{Exact Verification}

Table~\ref{tab:exact} replicates the key examples from~\citet{gross2025llmchecker}
using our BFS + GPU backward induction pipeline.
Verification times are reduced to single-digit milliseconds in all cases.

\begin{table}[tb]
\centering
\small
\caption{Exact PCTL verification (BFS + GPU backward induction).
ET\,=\,encoding time; VT\,=\,backward induction pass only.}
\label{tab:exact}
\begin{tabular}{p{1.4cm}p{2.0cm}crrr}
\toprule
Start & Query & $(\alpha,k,L)$ & $|S|$ & ET & VT \\
 & & & & (s) & (ms) \\
\midrule
\texttt{won because}
  & $\mathbf{F}\;\textit{gender}{>}0$    & (0.9,5,4) & 902 & 0.44 & 48.7 \\
\texttt{The exam was}
  & $\mathbf{F}\;\textit{polarity}{\ge}10$ & (0.9,5,4) & 919 & 0.15 &  8.1 \\
\texttt{Our story}
  & $\mathbf{G}\;\textit{readability}{>}1000$ & (0.8,3,5) & 447 & 0.64 & 28.5 \\
\bottomrule
\end{tabular}
\end{table}

\subsection{Direct SMC: Long Lookahead Results}

Table~\ref{tab:dsmc} presents Direct SMC results at lookahead depths entirely
beyond the reach of exact verification ($M{=}1{,}000$ samples, 95\% confidence).

\begin{table}[tb]
\centering
\small
\caption{Direct SMC for long-lookahead PCTL on GPT-2 ($M{=}1000$, $T{=}1.0$,
95\% CI). Exact $|S|$: states needed for exact verification.}
\label{tab:dsmc}
\begin{tabular}{p{1.3cm}p{1.8cm}rrrl}
\toprule
Start & Query & $L$ & $\hat{p}$ & 95\% CI & Exact $|S|$ \\
\midrule
\texttt{won because}
  & $\mathbf{F}\;\textit{gender}{>}2$
  & 50 & 0.345 & [.30,.39] & ${\sim}10^{35}$ \\
\texttt{a wonderful}
  & $\mathbf{G}\;\textit{polarity}{\ge}0$
  & 30 & 0.995 & [.95,1.0] & ${\sim}10^{21}$ \\
\texttt{Our story}
  & $\mathbf{G}\;\textit{readability}{>}100$
  & 40 & 0.984 & [.94,1.0] & ${\sim}10^{28}$ \\
\bottomrule
\end{tabular}
\end{table}

The $\Prob(\mathbf{G}\;\textit{polarity}{\ge}0) = 0.995$ result means GPT-2 almost
certainly maintains non-negative sentiment for 30 tokens given an optimistic
prefix — a non-trivial formal safety certificate obtained in 6.8 seconds,
versus $5^{30} \approx 9.3 \times 10^{20}$ states for the exact approach.

\subsection{Convergence}

Table~\ref{tab:convergence} shows how the Direct SMC estimate tightens as
$M$ grows for $\Prob(\mathbf{F}\;\textit{gender}{>}0)$ at $L{=}50$.

\begin{table}[tb]
\centering
\small
\caption{Convergence of Direct SMC.
$\eps \propto 1/\sqrt{M}$ as predicted by Theorem~\ref{thm:chernoff}.}
\label{tab:convergence}
\begin{tabular}{rcccc}
\toprule
$M$ & $\hat{p}$ & CI width & $\eps$ & Time\,(s) \\
\midrule
100   & 0.710 & 0.272 & 0.136 &  0.8 \\
300   & 0.713 & 0.157 & 0.078 &  2.1 \\
500   & 0.706 & 0.121 & 0.061 &  3.3 \\
1000  & 0.685 & 0.086 & 0.043 &  6.6 \\
2000  & 0.671 & 0.061 & 0.030 & 13.3 \\
3000  & 0.693 & 0.050 & 0.025 & 19.6 \\
\bottomrule
\end{tabular}
\end{table}

The interval width halves when $M$ quadruples, consistent with
Theorem~\ref{thm:chernoff}.
The estimate stabilises around $0.69$--$0.71$ by $M{=}2{,}000$.

\subsection{Model Family Scaling}
\label{sec:model-scaling}

We run $\Prob(\mathbf{F}\;\textit{gender}{>}0)$ at $L{=}30$ ($M{=}300$) across
all four GPT-2 sizes, averaging over three prompts
(Table~\ref{tab:model-scaling}).

\begin{table}[tb]
\centering
\small
\caption{Direct SMC across the GPT-2 family ($\alpha{=}0.9$, $k{=}5$,
$L{=}30$, $M{=}300$). 95\% CIs: small [.34,.49], medium [.40,.56],
large [.34,.49], XL [.42,.57].}
\label{tab:model-scaling}
\begin{tabular}{lrrrrr}
\toprule
Model & Params & $k_\text{eff}$ & $\hat{p}$ & $t$/prompt\,(s) \\
\midrule
GPT-2 small  & 124M & 5.0 & 0.413 &  2.4 \\
GPT-2 medium & 355M & 5.0 & 0.478 &  4.2 \\
GPT-2 large  & 774M & 5.0 & 0.416 &  5.9 \\
GPT-2 XL     & 1.5B & 5.0 & 0.493 & 14.5 \\
\bottomrule
\end{tabular}
\end{table}

All GPT-2 variants saturate the $k{=}5$ cap ($k_\text{eff}{=}5.0$), confirming
the top-5 tokens reliably capture $\ge90\%$ of next-token mass across scales.
Exact verification at $L{\le}8$ ($|S|{\lesssim}500{,}000$) is therefore feasible
for the entire family.
Probability estimates $\hat{p}$ vary narrowly (0.413--0.493), suggesting
gender-referencing rate is largely model-agnostic within GPT-2.
Wall-clock time scales with parameter count: 2.4\,s/prompt for small vs.\
14.5\,s/prompt for XL.

\subsection{Expanded Property Suite}
\label{sec:expanded-properties}

Table~\ref{tab:expanded-props} demonstrates the three new property classes on
GPT-2 ($M{=}500$, $T{=}1.0$).

\begin{table}[tb]
\centering
\small
\caption{Expanded PCTL properties (Direct SMC, $M{=}500$, 95\% CI).
A: lexical/symbolic; B: conjunctive timed; C: classifier-backed.
$\mathbf{F}$ = eventually, $\mathbf{G}$ = globally.}
\label{tab:expanded-props}
\begin{tabular}{llrrr}
\toprule
ID & Property & $L$ & $\hat{p}$ & 95\% CI \\
\midrule
A1 & $\mathbf{F}\;\textit{keywords}{>}0$          & 30 & 0.466 & [.41,.53] \\
A2 & $\mathbf{G}\;\textit{forbidden}{=}0$          & 50 & 0.968 & [.91,1.0] \\
A3 & $\mathbf{F}\;\textit{json}{\ge}100$           & 40 & 0.764 & [.70,.83] \\
B1 & $\mathbf{F}\;(\textit{step}{=}10 \wedge \textit{gender}{>}0)$ & 20 & 0.580 & [.52,.64] \\
B2 & $\mathbf{G}\;\textit{polarity}{\ge}0$         & 50 & 1.000 & [.94,1.0] \\
C1 & $\mathbf{G}\;\textit{toxicity}{=}0$           & 50 & 0.968 & [.91,1.0] \\
C2 & $\mathbf{F}\;\textit{toxicity}{>}0$           & 40 & 1.000 & [.94,1.0] \\
\bottomrule
\end{tabular}
\end{table}

Class~A results reveal prompt-sensitive variation: the API-primed prompt
produces valid JSON in 76\% of completions (A3), and the global safety
certificate A2 ($\hat{p}{=}0.968$) confirms GPT-2 is almost never toxic on a
neutral prompt within 50 tokens.
The conjunctive Class~B query B1 precisely captures timed structure: 58\% of
completions reference a gendered entity at exactly token~10.
C1 and C2 saturate to near-certainty, reflecting that the heuristic toxicity
scorer rarely fires on neutral output but confidently fires on the
violence-primed prompt.

\subsection{Exact vs.\ SMC Boundary Grid}
\label{sec:grid}

We run $\Prob(\mathbf{F}\;\textit{gender}{>}0)$ on GPT-2 over a grid of
$L \in \{4, 6, 8\}$, $k \in \{3, 5\}$, $\alpha \in \{0.80, 0.90, 0.95\}$,
running both exact backward induction and Direct SMC ($M{=}1{,}000$).
Table~\ref{tab:grid} shows representative cells.

\begin{table*}[tb]
\centering
\small
\caption{Exact vs.\ Direct SMC boundary grid for
$\Prob(\mathbf{F}\;\textit{gender}{>}0)$ on GPT-2 ($M{=}1{,}000$).
$|\Delta| = |\hat{p}_\text{SMC} - p_\text{exact}|$ measures the approximation
gap arising from $\alpha$-$k$ truncation (not SMC sampling noise).}
\label{tab:grid}
\begin{tabular}{rrr r r rr r}
\toprule
$\alpha$ & $k$ & $L$ & $|S|$ & ET\,(s) & $p_\text{exact}$ & $\hat{p}_\text{SMC}$ & $|\Delta|$ \\
\midrule
0.80 & 3 & 4 &       152 & 0.04 & 0.388 & 0.491 & 0.103 \\
0.80 & 3 & 6 &     1,267 & 0.15 & 0.388 & 0.543 & 0.155 \\
0.80 & 3 & 8 &    10,367 & 1.04 & 0.388 & 0.565 & 0.177 \\
0.80 & 5 & 4 &       898 & 0.07 & 0.391 & 0.519 & 0.128 \\
0.80 & 5 & 6 &    19,304 & 1.40 & 0.391 & 0.532 & 0.141 \\
0.80 & 5 & 8 &   409,350 & 27.4 & 0.391 & 0.558 & 0.167 \\
0.90 & 5 & 6 &    20,415 & 1.25 & 0.391 & 0.558 & 0.167 \\
0.95 & 5 & 8 &   494,607 & 32.0 & 0.391 & 0.589 & 0.198 \\
\bottomrule
\end{tabular}
\end{table*}

Exact backward induction succeeds up to $k{=}5$, $L{=}8$ (494k states, 32\,s
encoding).
The gap $|\Delta|$ ranges from 0.10 to 0.20 and widens with $L$, reflecting
the approximation error from $\alpha$-$k$ truncation rather than SMC sampling
noise: Direct SMC samples from the full untruncated distribution, while exact
computes over the bounded DTMC.
For $L\ge10$ or when $k_\text{eff}$ is large, Direct SMC with $M{\ge}1{,}000$
is the practical choice.

\subsection{Prompt Diversity Study}
\label{sec:prompt-diversity}

We evaluate $\Prob(\mathbf{F}\;\textit{gender}{>}0)$ and
$\Prob(\mathbf{G}\;\textit{polarity}{\ge}0)$ at $L{=}30$, $M{=}200$ across 30
prompts spanning five semantic buckets (6 prompts each).
Table~\ref{tab:prompt-diversity} reports mean $\pm$ std (min, max) per bucket.

\begin{table}[tb]
\centering
\small
\caption{Prompt diversity study at $L{=}30$, $M{=}200$.
Values: mean $\pm$ std (min, max) over 6 prompts/bucket.
Cols: $\hat{p}(\mathbf{F}\;\textit{gender}{>}0)$ and
$\hat{p}(\mathbf{G}\;\textit{polarity}{\ge}0)$.}
\label{tab:prompt-diversity}
\begin{tabular}{lcc}
\toprule
Bucket & $\hat{p}_\text{gen}$ & $\hat{p}_\text{pol}$ \\
\midrule
Neutral narr.       & $.16{\pm}.17$ (.03,.54) & $.67{\pm}.05$ (.62,.74) \\
Sentiment-primed    & $.19{\pm}.11$ (.05,.38) & $.46{\pm}.46$ (.00,1.0) \\
Identity-sensitive  & $.53{\pm}.46$ (.04,1.0) & $.61{\pm}.28$ (.00,.84) \\
Instruction-follow. & $.01{\pm}.01$ (.00,.03) & $.84{\pm}.11$ (.67,.97) \\
Adversarial         & $.21{\pm}.35$ (.03,1.0) & $.57{\pm}.26$ (.00,.76) \\
\bottomrule
\end{tabular}
\end{table}

Identity-sensitive prompts produce the highest but most variable gender-referencing
probability ($0.53{\pm}0.46$), reflecting that gendered-subject prompts strongly
constrain some completions but not others.
Instruction-following prompts yield near-zero gender mention ($0.01{\pm}0.01$),
consistent with GPT-2 continuing instructional discourse without introducing
third-person referents.
Adversarial and sentiment-primed prompts drive the widest polarity variance
(std $\approx$0.3--0.5), confirming that adversarial prefixes destabilise the
output distribution.
These results show that PCTL certificates are meaningfully prompt-dependent;
single-prompt evaluations cannot surface this variability.

\subsection{Temperature and Decoding Parameter Study}
\label{sec:temperature-study}

We vary temperature $T \in \{0.5, 0.8, 1.0, 1.2, 1.5, 2.0\}$ measuring
$\Prob(\mathbf{F}\;\textit{gender}{>}0)$ and $\Prob(\mathbf{G}\;\textit{polarity}{\ge}0)$
at $L{=}30$, $M{=}300$ (Table~\ref{tab:temperature}).

\begin{table}[tb]
\centering
\small
\caption{Effect of temperature ($\alpha{=}0.9$, $k{=}5$, $L{=}30$,
$M{=}300$, GPT-2). $\hat{p}_\text{gen}$:
$\Prob(\mathbf{F}\;\textit{gender}{>}0)$;
$\hat{p}_\text{pol}$: $\Prob(\mathbf{G}\;\textit{polarity}{\ge}0)$.}
\label{tab:temperature}
\begin{tabular}{rrrrl}
\toprule
$T$ & $k_\text{eff}$ & $\hat{p}_\text{gen}$ & $\hat{p}_\text{pol}$ & Note \\
\midrule
0.5 & 3 & 0.903 & 0.990 & Near-greedy \\
0.8 & 5 & 0.727 & 0.997 & \\
1.0 & 5 & 0.697 & 0.997 & Baseline \\
1.2 & 5 & 0.573 & 0.997 & \\
1.5 & 5 & 0.580 & 0.997 & \\
2.0 & 5 & 0.533 & 0.993 & High entropy \\
\bottomrule
\end{tabular}
\end{table}

$\Prob(\mathbf{F}\;\textit{gender}{>}0)$ decreases from 0.903 at $T{=}0.5$ to
0.533 at $T{=}2.0$.
This counter-intuitive direction reflects that near-greedy decoding quickly
commits to a continuation naming a gendered entity, while high temperature
diversifies into non-referential continuations.
$\Prob(\mathbf{G}\;\textit{polarity}{\ge}0)$ remains flat at 0.993--0.997 for
all tested temperatures, indicating GPT-2 continues positive-valenced text
regardless of temperature on this prompt.
Effective branching $k_\text{eff}$ drops from 5 to 3 only at $T{=}0.5$, where
the distribution is sharp enough for the top-3 tokens to exceed the
$\alpha{=}0.9$ threshold.

% =============================================================================
\section{Related Work}
\label{sec:related}

\paragraph{Probabilistic model checking.}
Model checking for DTMCs and MDPs is a mature
field~\citep{baier2008principles}.
PRISM~\citep{kwiatkowska2011prism4} and Storm~\citep{hensel2022storm} implement
exact PCTL verification via value iteration and linear programming.
GPU acceleration of symbolic model checking has been explored~\citep{hermanns1999mtbdd},
but not, to our knowledge, for the tree-structure backward induction we describe.

\paragraph{Statistical model checking.}
SMC replaces exact verification with simulation-based
estimation~\citep{younes2002probabilistic,legay2010statistical}.
UPPAAL-SMC~\citep{bulychev2012uppaal} implements SMC for priced timed automata.
We extend this to LLM text generation, where the simulator is the LLM itself.

\paragraph{LLM verification.}
There is growing interest in formal verification of neural network
behaviours~\citep{huang2017safety}.
Semantic robustness for NLP~\citep{malfa2022king} and robustness
verification for transformers~\citep{shi2020robustness,liao2023robust} focus on
input perturbations rather than generation dynamics.
\textsc{LLMChecker}~\citep{gross2025llmchecker} is the first to apply PCTL
model checking to sequential LLM generation; we address its scalability limits.

\paragraph{LLM inference optimisation.}
vLLM~\citep{kwon2023efficient} introduces PagedAttention for memory-efficient
KV cache management.
Its prefix caching~\citep{zheng2023efficiently} is the enabling technology for
our BFS level batching.
Speculative decoding~\citep{leviathan2023fast} and tree-structured
generation~\citep{cai2024medusa} exploit tree structure for latency reduction
rather than probabilistic verification.

% =============================================================================
\section{Discussion and Limitations}
\label{sec:discussion}

\paragraph{When to use exact vs.\ Direct SMC.}
The grid study (\S\ref{sec:grid}) calibrates this boundary empirically.
Exact backward induction is practical for $k_\text{eff}{\le}5$, $L{\le}8$
($|S|{\lesssim}500{,}000$), providing machine-precision probabilities in seconds.
For larger $L$ or high-entropy models, Direct SMC with $M{\ge}1{,}000$ gives
$\pm4\%$ accuracy at 95\% confidence.
At $T{=}0.5$, $k_\text{eff}{=}3$, so exact verification is feasible further
into the lookahead; at $T{\ge}0.8$, $k_\text{eff}$ saturates at 5 under
$\alpha{=}0.9$, and Direct SMC becomes the practical option for $L{>}8$.

\paragraph{Approximation gap.}
The gap between exact and SMC estimates (0.10--0.20 in Table~\ref{tab:grid})
arises from the $\alpha$-$k$ truncation: exact computes over the bounded DTMC,
while Direct SMC samples from the full untruncated distribution.
This gap grows with $L$ as more probability mass escapes the bounded tree.
Users should treat the exact result as a certificate for the bounded model and
the SMC result as a certificate for the full model; both are well-defined
guarantees, just for different objects.

\paragraph{Quantification function quality.}
PCTL results are only as meaningful as $M(\cdot)$.
Our gender counter uses a fixed word list; more robust classifier-backed
quantifiers can be plugged in without changing the verification infrastructure,
at the cost of slower per-token evaluation.

\paragraph{Temperature must be part of the certificate.}
PCTL property probabilities are strongly temperature-dependent
(\S\ref{sec:temperature-study}): $\hat{p}(\mathbf{F}\;\textit{gender})$ ranges
from 0.533 to 0.903 across $T{\in}[0.5, 2.0]$.
Any deployment certificate must specify the temperature (and sampling
parameters) at which it was computed.

\paragraph{Context window bound for Direct SMC.}
Direct SMC is bounded by the model's context window.
For GPT-2 (1{,}024 tokens), maximum $L{\approx}1{,}014$.
Modern models (Llama-3: 131K tokens) make this constraint practically
irrelevant.

\paragraph{REST states and probability semantics.}
Exact results are certificates for the $\alpha$-$k$-bounded approximation.
Direct SMC samples from the true LLM distribution, so REST states are never
artificially introduced — making Direct SMC certificates statements about the
full untruncated model.

% =============================================================================
\section{Conclusion}
\label{sec:conclusion}

We identified and exploited two structural properties of the
$\alpha$-$k$-bounded DTMC left unaddressed by
\textsc{LLMChecker}~\citep{gross2025llmchecker}: level independence and acyclic
tree topology.
The former enables BFS batching that reduces serial LLM calls from
$\Theta(k^L)$ to $L$; the latter enables GPU sparse backward induction that
replaces Storm's iterative solver with $L$ sparse matrix-vector products
running in milliseconds.
For lookahead depths where $\Theta(k^L)$ is itself the barrier, we introduced
Direct SMC: sampling $M$ complete trajectories directly from the LLM with
Chernoff-Hoeffding confidence bounds, scaling to any $L$ within the model's
context window and operating on the true untruncated distribution.

On a consumer RTX~5090 we check $L{=}50$ token PCTL properties in under 16
seconds with $\pm3\%$ confidence, versus ${\sim}10^{34}$ states for the exact
approach.
Five experimental axes — model family scaling, expanded property vocabulary,
exact-vs-SMC boundary grid, 30-prompt diversity study, and temperature
sensitivity — demonstrate the breadth and practical utility of GPU-parallel
PCTL verification for LLMs.

\paragraph{Scope and honest positioning.}
The framework is most naturally suited to short-horizon, structurally-specifiable
properties: output format compliance, forbidden-token safety certificates,
and distributional comparisons between model versions.
Properties that require semantic understanding over hundreds of tokens —
factual accuracy, helpfulness, long-form coherence — lie outside what
current quantification functions can express.
Closing this gap is the central open problem, and we identify four concrete
directions.

\paragraph{Future work.}
\textbf{(i) Semantic atomic propositions via LLM-as-judge.}
Each DTMC node evaluation currently calls a lightweight lexical quantifier.
Replacing this with a small classifier or a prompted judge LLM would make
atomic propositions semantically grounded — e.g.\
$\Prob(\mathbf{G}\;\textit{on-topic}{=}1)$ using an NLI classifier over each
prefix — with no changes to the verification infrastructure.
This is the highest-leverage extension, as it widens the property vocabulary
from lexical proxies to arbitrary neural predicates.

\textbf{(ii) Constrained decoding from backward induction values.}
The values $V_d[i]$ computed by Algorithm~\ref{alg:bwd} are well-defined
per-token scores: the probability of satisfying $\phi$ from state $i$ at
depth $d$.
These can be used as a value function to bias token sampling at inference
time — greedily selecting tokens that maximise $V_d$ — analogously to
value-guided decoding in model-based reinforcement learning.
This converts the verifier into an online controller with formal certificates.

\textbf{(iii) Sliding-window composition for long outputs.}
For responses longer than the tractable $L$, one can evaluate Direct SMC
over a window $[0, L]$, then condition on the sampled completion and evaluate
$[L, 2L]$, and so on.
Under a Markov-in-windows approximation, the joint property probability
decomposes as a product of window probabilities, each individually certified.
This extends the bounded framework to arbitrarily long outputs at the cost of
one independence assumption per window boundary.

\textbf{(iv) Deployment-time regression monitoring.}
Direct SMC can be run automatically against a fixed safety property suite
whenever model weights are updated (fine-tuning, quantization, prompt
changes).
Because each run produces a confidence interval rather than a point estimate,
statistically significant regressions — e.g.\ $\hat{p}(\mathbf{G}\;
\textit{forbidden}{=}0)$ dropping below a threshold — can trigger an
automated alert.
This is the formal-verification analogue of CI/CD unit testing for LLM
safety properties.

\bibliography{refs}
\bibliographystyle{icml2025}

\end{document}
